\caption{\small{\textbf{Overview of the probe generation pipelines.}
\textbf{(A) Factual Probe Generation}: We generate probes on a sentence-by-sentence basis to stay close to the source document and preserve reliable ground truth.
\underline{Preprocessing}: We first filter sentences that contain minimal knowledge, such as structural or discourse-only comments. We then use heuristics to remove sentences likely to yield low-quality probes, including sentences that are too short or dominated by \LaTeX{}, markup, or other formatting artifacts with no valid natural-language target for extraction.
\underline{Question Extraction}: From each filtered sentence, we extract 1--3 questions that capture the knowledge contained in the sentence.
\underline{Contextualization}: Questions are contextualized so they are clear, self-contained, and faithful to the knowledge being tested.
\underline{Cloze Conversion}: The questions are converted into cloze statements, placing the answer at the end.
\underline{Refinement}: We refine the statements, ensure mathematical content is written in \LaTeX{} when appropriate, and verify the correctness of the preceding steps.
\textbf{(B) Inference Probe Generation}: We generate inference probes to test conclusions supported by the document but not necessarily stated as isolated facts. We use a two-level approach. First, we break the document into sections and prompt an LLM to extract inference questions that require reasoning over the supporting text. This section-level pass encourages granular probes grounded in local evidence. Second, we repeat the step with the full document to allow more holistic probes that synthesize information across multiple facts or sections.
\underline{Cloze Conversion}: Similar to the factual process, these questions are converted into cloze statements with the answer at the end.
\underline{Refinement}: Statements are refined to ensure mathematical content is in \LaTeX{} where appropriate and to enforce proper self-containment by referencing the source document when needed.
\underline{Filtering}: A comprehensive filtering step asks an LLM to assess whether probes are too simple, imprecise, ambiguous, or confusing. This step removes roughly 15\% of probes on average.}}
